\documentclass[10pt,a4paper]{amsart}
\usepackage[margin=19mm]{geometry}
\usepackage{amsmath,amssymb,mathtools}
\usepackage[scr=rsfs,cal=euler]{mathalfa}
\usepackage{xfrac}
\usepackage[unicode,hidelinks,bookmarks=false]{hyperref}
\newcommand{\ie}{i.e.\ }
\newcommand{\cf}{cf.\ }
\newcommand{\resp}{respectively\ }
\newcommand{\Subsection}{\subsection}
\providecommand{\nobreakdash}{\nobreak\hbox{-}}
\setlength{\emergencystretch}{2em}
\multlinegap0pt
\usepackage{bm}
\usepackage{bbm}
\usepackage{dsfont}
\usepackage{xspace}
\usepackage{cancel}
\usepackage{mathtools}
\usepackage{amsthm}
\makeatletter
\@ifundefined{lemma}{\newtheorem{lemma}{Lemma}}{}
\@ifundefined{theorem}{\newtheorem{theorem}{Theorem}}{}
\makeatother
\newcommand{\CS}{\textnormal{\textsf{CS}}}
\newcommand{\Mw}{\mathcal{M}}
\newcommand{\Prop}{\mathtt{Prop}}
\newcommand{\Ra}{\Rightarrow}
\newcommand{\R}{\mathcal{R}}
\newcommand{\Seq}{\mathfrak{S}}
\newcommand{\labCS}{{\textnormal{labCS}^\infty}}
\newcommand{\la}{\langle}
\newcommand{\ra}{\rangle}
\newcommand{\spa}{\;|\;}
\newcommand{\tree}{\mathcal{T}}
\title[A Non-Wellfounded and Labelled Sequent Calculus for Bimodal ]{A Non-Wellfounded and Labelled Sequent Calculus for Bimodal Provability Logic}
\author{Justus Becker}
\date{Source: arXiv:2506.14307v1 (CC BY 4.0)}
\begin{document}
\maketitle

\noindent\emph{Source and licence.} A Non-Wellfounded and Labelled Sequent Calculus for Bimodal Provability Logic. Version arXiv:2506.14307v1 (CC BY 4.0), distributed under the Creative Commons Attribution 4.0 International licence, as recorded in the arXiv licence metadata for this exact version and shown on the abstract page https://arxiv.org/abs/2506.14307v1. Full rights notice: licenses/arxiv-2506.14307-CC-BY-4.0.txt.\par\smallskip

\noindent\textit{Editorial scope.} Counted unit: the Theorem whose source label is \texttt{thm:complCS} in the article (source0.tex lines 568--570, source label \texttt{thm:complCS}), delivered together with the article's own complete original proof of it (source0.tex lines 571--606, one \texttt{proof} environment of 36 lines). No mathematics is written, repaired or re-derived: statement and proof are the article's own text line for line. Required context retained uncounted: lem:FalsityCS (lines 430--436); lem:fulsat (lines 537--539). Every definition, display and label of those lines is retained unchanged. Omitted from this package and not redistributed: 1--429, 437--536, 540--567, 607--635. Nothing inside a retained excerpt is omitted or altered: every retained body is delivered exactly as the article's lines are written, so no omission receipt of any kind accompanies this package. Citations: the retained excerpts carry no citation command at all, so no bibliography entry of the article is transcribed and no reference is redistributed. Reference adaptation: none was needed and none was made; every reference inside the delivered unit resolves inside it, so no unresolved reference or citation is produced. Printed slips kept literal: no printed word, symbol, hypothesis or number of the counted statement or of its proof was altered. Layout and class: the article's own class is \texttt{llncs} with packages including amsmath, hyperref, tikz-cd, tikz, amssymb, graphicx, fontenc, ebproof, mdframed, amsthm, stmaryrd, adjustbox, soul, xcolor; the wrapper is the lane's pinned \texttt{amsart} editorial wrapper at 10pt on A4, which restates the theorem-like environments the retained text needs and adds the article's own macro definitions that the retained text needs (\texttt{\textbackslash CS}, \texttt{\textbackslash Mw}, \texttt{\textbackslash Prop}, \texttt{\textbackslash R}, \texttt{\textbackslash Ra}, \texttt{\textbackslash Seq}, \texttt{\textbackslash la}, \texttt{\textbackslash labCS}, \texttt{\textbackslash ra}, \texttt{\textbackslash spa}, \texttt{\textbackslash tree}), with extra wrapper packages (bm, bbm, dsfont, xspace, cancel, mathtools). The delivered numbering is the unit's own. Assets: no retained span contains a figure environment, a graphics-inclusion command, a PDF-page inclusion, a table or a TikZ picture; no image file is copied into this package, the package declares no asset dependency and no image file is redistributed with it. Licence: version arXiv:2506.14307v1 of this article is distributed under the Creative Commons Attribution 4.0 International licence, as recorded in the arXiv licence metadata for this exact version and shown on the abstract page https://arxiv.org/abs/2506.14307v1. A full rights notice is delivered at licenses/arxiv-2506.14307-CC-BY-4.0.txt. Characters: every retained excerpt is pure ASCII, so the delivered engine, XeLaTeX with its default Latin Modern text font, reports no missing character.
\begin{lemma}[Local Soundness] \label{lem:FalsityCS}
    Let $\Mw$ be a Carlson model and let
    \[
    \frac{\Seq_1 ... \Seq_n}{\Seq}
    \]
    be a rule of $\labCS$ with at least one premiss, then: $\Mw \Vdash \Seq$ iff $\Mw \Vdash \Seq_i$ for all $1 \leq i \leq n$.
\end{lemma}

\begin{lemma} \label{lem:fulsat}
    If $\Seq$ is a fully saturated sequent, then there is a countermodel ${\Mw \nVdash \Seq}$. 
\end{lemma}

\begin{theorem}[Completeness of $\labCS$] \label{thm:complCS}
    For any sequent $\Seq$: If $\CS \Vdash \Seq$ then $\labCS$ $\vdash \Seq$. Furthermore, if a sequent is not provable in $\labCS$, there is a countermodel $\Mw \nVdash \Seq$.
\end{theorem}

\begin{proof}    
    We prove the contrapositive.
    Assume $\labCS \nvdash \Seq$. Perform the proof search described above on $\Seq$ to obtain a derivation $\tree$. Consider some branch $(\Seq_i)_{i < \omega}$ in $\tree$ such that it either ends with a fully saturated sequent, or it is infinite but not progressing. Such a branch must exist by $\labCS \nvdash \Seq$ as we otherwise would have shown the opposite. 
    Write $\Seq_i = \R_i , \Gamma_i \Ra \Omega_i$ and define a sequence of models $(\Mw_i)_{i < \omega}$ with $\Mw_i = \la W_i , M_{0 \, i}, M_{1 \, i}, \prec_i, V_i \ra$ as follows.
    \begin{itemize}
        \item[] $W_i := Lab(\Seq_i)$
        \item[] $M_{0 \, i} := \{ x \in W_i \spa \text{there is some } y \text{ such that } yRx \in \R_i \}$
        \item[] $M_{1 \, i}:= \{ x \in W_i \spa \text{there is some } y \text{ such that } ySx \in \R_i \}$
        \item[] ${\prec_i} := \{ (x,y) \in W_i \times W_i \spa xRy \in \R_i \text{ or } xSy \in \R_i \}$
        \item[] $V_i(p) := \{ x \in W_i \spa x : p \in \Gamma \}$
    \end{itemize}
    
    If $(\Seq_i)_{i < \omega}$ is finite and ends with a fully saturated sequent $\Seq_n$, by lemma \ref{lem:fulsat} $\Mw_n \nVdash \Seq_n$. Then, by lemma \ref{lem:FalsityCS} also $\Mw_n \nVdash \Seq$.
    
    If $(\Seq_n)_{n < \omega}$ is infinite, then any potential trace in $(\Seq_n)_{n < \omega}$ must be eventually constant. Consider the limit model $\Mw = \la W , \prec , M_0 , M_1 , V \ra$ with $W : =\bigcup_{i < \omega} W_i$, ${\prec} : =\bigcup_{i < \omega} \prec_i$, $M_0 : =\bigcup_{i < \omega} M_{0 \, i}$, $M_1 : =\bigcup_{i < \omega} M_{1 \, i}$, and $V(p) : =\bigcup_{i < \omega} V_i (p)$. 
    By every trace of $(\Seq_n)_{n < \omega}$ being eventually constant, $\prec$ must be conversely wellfounded, otherwise an infinite $\prec$-chain would allow to obtain a progressing trace. (We do not consider the case for $\prec$ having a cycle as it must form a tree if the root of the derivation is cycle free.) Therefore, $\Mw$ must be a Carlson model. We now show that for any $\Seq_i$ in the branch we have $\Mw \nVdash \Seq_i$. We do this similarly to the proof of lemma \ref{lem:fulsat} by induction over the complexity of formulas occurring in $\Seq_i$, for all $\Seq_i$ simultaneously. 
    
    Consider any sequent $\Seq_n$ in the branch and a formula $A$ such that $x:A$  occurs in $\Seq_n$. It then suffices to show that $\Mw, x \Vdash A$ if $x:A \in \Gamma_n$ and $\Mw , y \nVdash B$ if $x:A \in \Omega_n$.

    $A \in \Prop$: If $x:A \in \Gamma $, the claim follows by definition. If $x:A \in \Omega$, it follows from the fact that the branch did not close during the proof strategy, thus Id cannot be applied.

    $A=\bot$: We know that $x:A \in \Gamma$ cannot happen because then the rule $\bot$ could have been applied, making the branch finite in the proof search strategy. If $x:A \in \Omega$, the formula is falsified by definition.

    \sloppy
    $A= A_1 \to A_2$: By the strategy $A$ must be eventually principal. 
    If ${x:A \in \Gamma}$, there is a sequent $\Seq_m$ with ${n < m}$ such that either ${x: A_1 \in \Omega_m}$ or ${x:A_2 \in \Gamma_m}$.
    By induction hypothesis $\Mw , x \nVdash A_1$ or $\Mw , x \Vdash A_2$, thus in both cases ${\Mw , x\Vdash A_1 \to A_2}$. 
    If $x:A \in \Omega$, there is a sequent $\Seq_m$ with $n < m$ such that $A_1 \in \Gamma_m$ and $x: A_2 \in \Omega_m$. By induction hypothesis $\Mw , y \Vdash A_1$ and $\Mw ,y \nVdash A_2$, so $\Mw ,y \nVdash A_1 \to A_2$.
    
    $A = \Box A_1$: If $x:A \in \Gamma$, there must be some $\Seq_m$ with $n < m$ such that $z : A_1 \in \Gamma_m$ for any $z$ with $xR z \in \R_m$. $\Mw, x \Vdash \Box A_1$ holds vacuously if no such $z$ exists. By induction hypothesis, $\Mw , z \Vdash A_1$ and by the definition of $\prec$ and of the set $M_0$: $\Mw , x \Vdash \Box A_1$.
    If $x:A \in \Omega$, there is some $\Seq_m$ with $n < m$ such that $z : A_1 \in \Omega_m$ and $xRz \in \R_m$. By induction hypothesis $\Mw , z \nVdash A_1$, $y \prec z$, and $z \in M_0$. Therefore, $\Mw , y \nVdash \Box A_1$.
    
    The case of $A = \triangle A_1$ works similarly to the previous one. 

    As $\Mw \nVdash \Seq_i$ for all $i < \omega$, we gain as a corollary $\Mw \nVdash \Seq$.  
\end{proof}

\end{document}
